\section{Certificate specifications}
\label{supp:templates}

This section prints the certificate specifications and definitions of the catalog that the main
paper summarizes in \Cref{tab:templates}; E3 and E13 are printed in the main paper
(\Cref{tmpl:e3,tmpl:e13}). Each Template gives its data, the certificate that carries the difficult
content, and the conclusion that Lean derives from a complete certificate. For a conditional
classification the certificate is a complete status record: one carried, branch-valid status record
for the claim, and the same status on every genuine record for that claim. After each Template, a
\emph{Meaning} sentence says what the conclusion amounts to, and an \emph{Example} gives one status
realized in the finite computations of \Cref{supp:evidence}; the examples are finite checks of the
classifiers, not instances of the entries.

\subsection{Endogeny}\label{sec:tmpl-endogeny}

\begin{schema}[E1: internalization]\label{tmpl:e1}
\emph{Data.} An E-system $\mathbb S$, a challenge history $H$, a recurring challenge class
$\mathcal C$ with horizon, and a viability readout $v$ with $\Delta(Q_t, F_{\mathcal C}, v) \ne
\varnothing$ at challenged times.
\emph{Certificate.} A complete internalization status record with five statuses: slack, externally
subsidized, endogenously repairing, stressed and collapsing. The challenge binds. The set
$\Delta_{\mathrm{endo}}$ is empty; it collects the discharging entries that have a forbidden source,
lack generation, scope or carriedness, use an off-kernel move, or omit their audit. The certified
inputs listed below are supplied.
\emph{Conclusion.} The five statuses are mutually exclusive, and the status is endogenously
repairing. The repair family contained in that branch has, at each challenged episode:
\begin{itemize}
  \item a carried defect and a generated move;
  \item the installed join $Q_{t+1} = Q_t \vee R_t$;
  \item a strict reduction
    $\Delta(Q_{t+1}, F_{\mathcal C}, v) \subsetneq \Delta(Q_t, F_{\mathcal C}, v)$ on the challenged
    fibres;
  \item a committed-state source, a kernel transition, gate passage and re-audit.
\end{itemize}
\end{schema}

\emph{Meaning.} If the record says the system repairs endogenously and its competitors are ruled out,
then every credited repair is a carried join that strictly shrinks the obstruction; the repair family
itself is a field of the certificate, not constructed from recurrence. The certified inputs are
no-free-distinction \citep{Tsiokos2026FoundIV}, promotion-gate soundness \citep{Tsiokos2026FoundIII},
a bridge from measured discharge to the absence of residual accrual, and family-level budget
feasibility. Perpetual self-extension under a challenge that is always new is the law of
\Cref{sec:e1-forcing}. A further statement has the same conditional form: if no repair generator is reachable, a supplied bridge excludes an
endogenous family and $\Delta_{\mathrm{endo}} = \varnothing$, then the status is slack, stressed or
collapsing, and a binding challenge excludes slack.
\emph{Example.} With its repair generator removed, a finite carrier under a split challenge is
classified stressed, stressed, collapsing at three successive times, with residuals $1, 2, 3$.

\begin{schema}[E4: repair compilation]\label{tmpl:e4}
\emph{Data.} A repair $R$ invoked at least twice for a challenge class $\mathcal C$, and a candidate
lower operator or support gate (sort P1 or P2) that would replace the invocation.
\emph{Certificate.} A compilation case: the candidate is carried; a certified top-down channel links
it to at least two admissible interventions with matched controls; the promotion record is accepted;
a memory-only comparator fails to reproduce the descent; the repair is not same-family saturated; an
annotation record permits decompilation; in-class descent is verified,
$\Delta(Q^{\mathrm{compiled}}, F_{\mathcal C}, r) = \varnothing$; the higher invocation is silent;
and the higher-priority statuses are excluded.
\emph{Conclusion.} The status is compiled; the five statuses exposed, compiling, compiled,
mis-compiled and decompiled are mutually exclusive.
\end{schema}

\emph{Meaning.} A repeated repair counts as moved into lower support only when every gate passes;
the classification then separates a successful compilation from one that fails outside its class
(mis-compiled) or has been undone (decompiled).
\emph{Example.} A candidate whose top-down channel is blocked is classified compiling; the same
candidate with all gates passing is compiled.

\begin{definition}[E5: reclosure collapse]\label{def:e5}
Fix a challenge class, a horizon and a complete inventory of typed rescue candidates --- carried
repairs, boundary updates, apparatus reclosures and probe acquisitions --- each with its own move
record, budget witness, kernel transition and viability-descent certificate. The system is
\emph{collapsed} if no candidate is both admissible within budget and viability-descending. It is
\emph{irreversibly collapsed} if moreover the repair generator is unreachable and the apparatus-level
viability kernel is empty, both certified for the same system, challenge, horizon and scope.
\end{definition}

\begin{lemma}[collapse consistency]\label{lem:e5}
An admissible in-budget descending candidate refutes collapse. Under a complete status record, the
seven statuses viable, stressed, subsidized, suspended, collapsed-recoverable,
collapsed-irreversible and revived are mutually exclusive. After a subsidy is withdrawn, a certified
absence of active subsidy and the complete record classify the system without presupposing collapse.
\end{lemma}

\emph{Meaning.} Collapse is the absence, in the declared inventory, of any admissible in-budget rescue
that restores viability descent;
suspension (intact apparatus, operations gated off) and revival (a later rescue with a new apparatus
lineage) are separate statuses. The first claim of the lemma holds by definition.
\emph{Example.} A carrier with an active kernel but no rescue move is collapsed-recoverable.

\subsection{Priced access and control}\label{sec:tmpl-access}

\begin{schema}[E7: alarm]\label{tmpl:e7}
\emph{Data.} A blind-spot witness $z$ with $z^*(\Xi_C(D \mid L) - \Omega)z > 0$, coupled within the
horizon to a failure of a declared viability probe.
\emph{Certificate.} A unique total disposition package: one carried, branch-valid disposition record,
and the same kind on every genuine record at this witness and horizon.
\emph{Conclusion.} Exactly one disposition holds:
\begin{itemize}
  \item \emph{currentization}: a carried rewrite of the allocator's constraint set that admits the
    witness and is charged to a forced-exposure ledger entry;
  \item \emph{boundary nonclosure}: a failed viability descent with a carried status;
  \item \emph{lawful discount}: for protocol artifact, currentizable residue, out-of-horizon scope or
    dissolved coupling.
\end{itemize}
A currentization changes the constraint set.
\end{schema}

\emph{Meaning.} To \emph{currentize} a blind-spot direction is to make it part of what the system
currently attends to, by changing the constraints of its allocation. A dangerous blind spot is then
either currentized in this way, declared an unclosed boundary, or discounted for a stated reason. A
move that wins the allocation of \Cref{sec:allocation} inside the old constraint set is attention,
not alarm. A related statement, \emph{alarm fatigue} (E7.1), is argued on paper only: if
viability-coupled excess persists and demand exceeds a finite carried preemption capacity, then the
channel saturates, a carried threshold repair is charged with the original lineage kept, or the
excess stays statused.
\emph{Example.} Protocol-artifact witnesses at thresholds $0$ and $1$ are lawfully discounted.

\begin{definition}[E8: control price]\label{def:e8}
A carried signal is a \emph{control-price component} for constraint $j$ of a maximization problem
with optimal value $\Phi^*(b)$ under constraints of level $b = (b_j)$ when its readout is certified, it is linked to the ledger line of $j$, its held-out
predictions and its stability under budget perturbation pass, and its value matches the multiplier
$\lambda_j$ of a linked multiplier witness when $j$ binds and vanishes when $j$ is slack. A
\emph{field} is a complete inventory of components; a compressed summary
$a = \mathcal C(\lambda_1, \ldots, \lambda_k)$ is \emph{lawful} when it predicts admissibility or
policy changes beyond the external quotient and collapses when all its components are slack.
\end{definition}

\begin{lemma}[control-price consistency]\label{lem:e8}
A component-lawful case yields its shadow-price component and slack collapse. Summary lawfulness is
equivalent to predictive legitimacy together with slack collapse, by definition. Under a complete
status record the seven statuses slack-obstructed, proxy-obstructed, summary-redescription,
component-lawful, field-lawful, summary-lawful and incomplete-or-unpriced are mutually exclusive.
\end{lemma}

\emph{Meaning.} A control signal counts as a price only when it is linked to the ledger line of its
own constraint, matches that constraint's multiplier, passes its held-out and stability checks, and
vanishes when the constraint is slack. When $\Phi^*$ is differentiable and the envelope theorem
applies, $\lambda_j = \partial\Phi^*/\partial b_j$
\citep[\S 5.6.3]{BoydVandenberghe2004}; this reading is on paper. \emph{Audit capture} (E8.1) is the
record in which a budget-setting move raises the carrier's own residual budget, the blind spot
against the original lineage stays positive, and the ledger reports health.
\emph{Example.} A signal of $1/4$ on a slack constraint is slack-obstructed. A budget move from $10$
to $15$ with the original-lineage blind spot still positive and a healthy report is capture evidence.

\begin{schema}[E10: cognitive demarcation]\label{tmpl:e10}
\emph{Data.} A carried package $\Theta$ and a challenge class $\mathcal C$.
\emph{Certificate.} A complete cognitive status record over eight statuses: schedule trap,
scaffolding, structural path only, decodable correlate, repair without gate, gate without repair,
cognitive, non-cognitive.
\emph{Conclusion.} The statuses are mutually exclusive. A cognitive status yields
\begin{equation}\label{eq:e10}
  \mathsf{RepairDescent}(\Theta, \mathcal C) \wedge \mathsf{Gate}(\Theta, \mathcal C) \wedge
  \mathsf{Prov}(\Theta, \mathcal C),
\end{equation}
and conversely \eqref{eq:e10} together with the absence of the six higher-priority kinds of
evidence yields the cognitive status.
\end{schema}

\emph{Meaning.} A package with cognitive status repairs, gates action and is the system's own, and
no higher-priority evidence is present: $\mathsf{RepairDescent}$ is a strict reduction of an access or action
obstruction by a quotient the package installs; $\mathsf{Gate}$ is a certified top-down kernel gate
for two distinct carried interventions with matched controls and different induced kernel supports;
and $\mathsf{Prov}$ is carried formation and invocation of the package, linked to $\mathcal C$. The
same form gives \emph{intention} (E10.1: a carried restriction of action support toward a target,
strictly before the outcome, certified and coherent under every listed future probe) with five
statuses, and \emph{goal} (E10.2: a carried target quotient that persists as the criterion for at
least two distinct carried routes) with four.
\emph{Example.} A package that an observer can decode but that has no certified difference-making
gate is a decodable correlate.

\subsection{Stack, individuation and transport}\label{sec:tmpl-stack}

\begin{schema}[E11: institutional rewrite]\label{tmpl:e11}
\emph{Data.} A formed upper package $\Theta$ over a lower kernel $K_\ell$ with closure $E_\ell$, and two
distinct admissible upper interventions $\Theta_i \ne \Theta_j$ with matched lower controls.
\emph{Certificate.} A complete status record over inert, conditioning, stack-active and constitutive,
with structural evidence consisting of a certified top-down channel linking the interventions to
\[
  \supp K_\ell^{\Theta_i} \ne \supp K_\ell^{\Theta_j}
  \quad\text{or}\quad
  E_\ell^{\Theta_i} \ne E_\ell^{\Theta_j}.
\]
\emph{Conclusion.} The four statuses are mutually exclusive. Stack-active requires structural
evidence, an empty stack obstruction and no constitutive evidence; constitutive requires, in
addition, a separate compilation certificate. A change of probabilities on unchanged support is
conditioning, and a label whose effect washes out is inert.
\end{schema}

\emph{Meaning.} An upper layer rewrites a lower one only if it changes which lower transitions are
possible or how the lower layer completes, not merely how likely they are.
\emph{Example.} A weight change from $10^{-3}$ to $10^{-2}$ on an allowed transition is conditioning.

\begin{schema}[E12: individuation]\label{tmpl:e12}
\emph{Data.} An E-system and a declared list of candidate subcarriers.
\emph{Certificate.} For a candidate $I$: repair, budget and record closure on $I$,
\[
  \Repair(I) \subseteq I, \qquad \Budget(I) \subseteq I, \qquad \Record(I) \subseteq I,
\]
an interface quotient $B(I)$ certified as the coarsest one through which every declared viability
probe descends, and self-maintenance of $B(I)$ by $I$; together with maximality among candidates
and the absence of a distinct overlapping maximal candidate. A complete status record ranges over
integrated, federated, subsidiary, platform-dependent, shadow and non-individuated.
\emph{Conclusion.} The status is integrated; the six statuses are mutually exclusive.
\end{schema}

\emph{Meaning.} An individual is a maximal declared candidate, with no overlapping maximal competitor,
whose repairs, budget and records stay inside it and whose coarsest viability-sufficient boundary it
maintains itself.
\emph{Example.} A candidate whose repairs and records close but whose budget is paid from outside is
subsidiary.

\begin{proposition}[a singleton candidate]\label{prop:e12-singleton}
If the candidate list is $[I]$ and $I$ passes the repair, budget, record, viability-boundary and
self-maintenance tests, then $I$ is integrated: it is maximal, and no distinct candidate overlaps it.
\end{proposition}

\begin{proof}
The only candidate is $I$, so maximality and the absence of a competitor are immediate.
\end{proof}

The semantic form of E12 --- repair closure, boundary maintenance, maximality and the exclusion of
overlapping rivals, for a finite system --- is the law of \Cref{sec:e12-semantic}; the Template adds
budget and record closure and a boundary certified as the coarsest viability-sufficient one.
\emph{Operational selfhood} (E12.1) is argued on paper only: at the cognitive scale, a \emph{self}
is a carried common repair ledger through which the system's role, memory, action and viability
probes all descend.

\subsection{Memory and adaptive capacity}\label{sec:tmpl-memory}

\begin{schema}[E14: reconsolidation]\label{tmpl:e14}
\emph{Data.} A carried record $m_t$, a retrieval transport $T_c(m_t)$ in context $c$, and a retrieval
conflict: a split pair against the record's own declared claim after the transported record is
adjoined.
\emph{Certificate.} Complete inventories of lawful repair, strict coarsening and statused outcomes,
with provenance records for every mutation, and a complete status record over nine statuses:
provenance defect, silent rewrite, ordinary read, outcome collision, unrealized disposition, record
repaired, record coarsened, statused unresolved, unstatused conflict.
\emph{Conclusion.} Once the first five statuses and unstatused conflict are excluded, exactly one
lawful disposition holds:
\[
\begin{aligned}
  &\text{repair:} && m_{t+1} \text{ installs the record join of } m_t \text{ and } R_c;\\
  &\text{coarsening:} && q^- \preceq q^+ \text{ and } q^+(x) = q^+(x'),\ q^-(x) \ne q^-(x')
     \text{ for some } x, x';\\
  &\text{retained debt:} && m_{t+1} = m_t \text{ with a statused residual charged to } \Lambda_S.
\end{aligned}
\]
\end{schema}

\emph{Meaning.} When recalling a record exposes a conflict and the five defect statuses and unstatused
conflict are excluded,
the record is lawfully repaired, made strictly coarser, or kept with the conflict charged as debt; the
other statuses record defects in the retrieval or in its bookkeeping.
\emph{Example.} The same source record is repaired for one retrieval claim and kept as statused
unresolved debt for another.

\begin{schema}[E15: offline reclosure]\label{tmpl:e15}
\emph{Data.} Carried closure debt $D_t$ built only from linked statused E14 residuals, route residues
and adequacy residuals; online capacity $\kappa_{\mathrm{on}}$; debt accrual $a$; and an allocation
record as in \Cref{prop:e15-identity}.
\emph{Certificate.} A complete status record over eight statuses: deficit-online counterexample,
decorative offline, duty-cycle mismatch, skipped-offline cascade, alternation required, online
sufficient, offline optional, offline reclosure rejected.
\emph{Conclusion.} The statuses are mutually exclusive. The alternation-required case yields a strict
online deficit $a > \kappa_{\mathrm{on}}$, recurrent offline phases that gate exchange to zero and
assign the freed allocation to internal discharge with linked debt discharge in each phase, and an
observed duty cycle within a tolerance declared before the horizon.
\end{schema}

\emph{Meaning.} A system that accrues debt faster than it can discharge online is certified to need
alternating offline phases only when the phases, the freed budget and the duty cycle are all on
record; \Cref{prop:e15-identity} gives the capacity of such a phase.
\emph{Example.} At debt accrual $5$ the claim is alternation required; at accruals $7/2$ and $4$ it is
online sufficient.

\begin{schema}[E16: adaptability]\label{tmpl:e16}
\emph{Data.} Two admissible routes $\gamma, \eta$ on one comparison support that end at the same
current practice quotient; a held-out challenge family; complete endpoint-conditioned trial
populations and their repair-capacity distributions.
\emph{Certificate.} A complete status record over eight statuses: support confound, artifact, flat,
flattenable, currentizable slack, dissipative, coherent adaptability, adaptability rejected; with the
controls for artifact, flattening, currentization, continuation, support confounding and bounded
perturbation.
\emph{Conclusion.} The statuses are mutually exclusive. Coherent-adaptability evidence yields that the
two endpoints are equivalent for every current event, and two claim-linked probability records for
the same challenge class with different probabilities.
\end{schema}

\emph{Meaning.} Two histories that look the same now are recorded as leaving different future
repair capacity, after every declared control has failed to explain the difference away; the
difference is read from the record, not identified as an effect of the route.
\emph{Example.} Two routes to current-event-equivalent endpoints give repair probabilities $3/4$ and
$1/4$ for the same challenge; six control regimes realize one status each, among them coherent
adaptability.
